\begin{claim}{D.50}{Proved}{The finite two-subgroups of the full punctured-torus group}
Let
$T=\Out_{\rm lac}(\widehat F(x,y);
\overline{\langle[x,y]\rangle})$ be the full
once-punctured-torus loop-outer group.
The possible abstract isomorphism types of its finite
$2$-subgroups are exactly
\[
                  1,\quad C_2,\quad C_4,\quad
                         C_2\times C_2,\quad D_8,
                                                        \tag{D.50a}
\]
where $D_8$ has order eight. All five types occur.
For a subgroup of type $C_2\times C_2$ or $D_8$,
the full multiplier takes its values in the globally
constant signs $\{1,-1\}\subset\widehat\Z^\times$.
In $C_2\times C_2$ there is one positive involution
and two negative involutions. In $D_8$ the rotation
subgroup has multiplier $1$ and the four reflections
have multiplier $-1$.

For an arbitrary cyclic subgroup of order four,
a generator $a$ satisfies the already proved
restrictions
\[
           \lambda(a)^2=1,\qquad \lambda(a)_2=1,
           \qquad \lambda(a^2)=1.                       \tag{D.50b}
\]
The abstract classification (D.50a) does not establish
conjugacy of arbitrary such subgroups, or replace
(D.50b) by a global multiplier determination.
\end{claim}
\begin{proof}
D.33 makes the actual homology homomorphism
\[
                     A_2:T\longrightarrow\GL_2(\Z_2)
                                                        \tag{D.50c}
\]
injective on every finite $2$-subgroup $K$ and proves
$|K|\leq8$. We classify its image using that same
specified injection.

No matrix in $\GL_2(\Q_2)$ has order eight:
an eigenvalue of order eight has minimal polynomial
$X^4+1$ over $\Q_2$, which is irreducible by
Eisenstein after $X\mapsto X+1$.
The same argument excludes higher powers of two.
Every matrix $R$ of order four has minimal
polynomial $X^2+1$, since this polynomial is
irreducible over $\Q_2$ and $R^2\ne I$.
Consequently
\[
                         R^2=-I.                       \tag{D.50d}
\]
An exponent-two subgroup is abelian, because
$(ab)^2=a^2=b^2=1$ implies $ab=ba$.
Its matrices are commuting diagonalizable
involutions over $\Q_2$, and simultaneous
diagonalization embeds it in
$\{\operatorname{diag}(\epsilon_1,\epsilon_2):
\epsilon_i=\pm1\}$. Its order is at most four.
Thus an image of order eight contains an $R$
of order four.

The matrix centralizer of $R$ is the multiplicative
group of the field $\Q_2[R]\cong\Q_2(i)$.
For completeness, choose a nonzero vector $v$:
the vectors $v,Rv$ are independent, since an
eigenvalue in $\Q_2$ would contradict the
irreducibility of $X^2+1$.
A commuting endomorphism is determined by its
value on $v$ and is therefore $aI+bR$.
Every element of the finite $2$-subgroup in this
field has order at most four by the preceding
matrix argument. It is consequently one of
the at most four roots of $X^4-1$.
Since $\langle R\rangle$ already has four elements,
\[
                       C_K(R)=\langle R\rangle
                                                        \tag{D.50e}
\]
for a subgroup of order eight.

Choose $S\in K\setminus\langle R\rangle$.
The index-two subgroup $\langle R\rangle$ is normal;
by (D.50e) conjugation by $S$ induces its nonidentity
automorphism. Hence
\[
          SRS^{-1}=R^{-1},\qquad
                    S^2\in\{I,R^2\}.                  \tag{D.50f}
\]
The second inclusion follows because $S^2\in
\langle R\rangle$ and $S$ fixes its own square,
whereas it inverts $R$.
The alternative $S^2=I$ gives the dihedral
presentation of order eight. We exclude the
alternative $S^2=R^2=-I$ explicitly.

In the basis $(v,Rv)$ used above one has
\[
 R=\begin{pmatrix}0&-1\\1&0\end{pmatrix}.
\]
The equation $SR=-RS$ forces
\[
             S=\begin{pmatrix}a&b\\b&-a\end{pmatrix},
             \qquad S^2=(a^2+b^2)I.
                                                        \tag{D.50g}
\]
Thus the second alternative would require
$a^2+b^2=-1$ in $\Q_2$.
Here is the exact valuation obstruction.
Put $m=\min(v_2(a),v_2(b))$ and write
$a=2^m u$, $b=2^m v$, with $u,v\in\Z_2$
and at least one a unit. If both are units,
$u^2+v^2\equiv2\pmod8$, so this sum has
valuation one. If exactly one is a unit,
the sum has valuation zero.
On the other hand the equation says
\[
                      u^2+v^2=-2^{-2m};
                                                        \tag{D.50h}
\]
the valuation of its right side is even.
Both-unit values are therefore impossible.
The other case forces $m=0$; reduction modulo
four then gives $1=-1$, also impossible.
Zero entries cause no exception to this
argument. This excludes the quaternion
alternative and leaves only $D_8$ at order
eight. Orders four, two, and one give the
remaining groups in (D.50a).

We specify actual full-group realizations.
On $\widehat F(x,y)$ define
\[
 \begin{array}{c|cc}
          &x&y\\ \hline
 r        &y&x^{-1}\\
 s        &y&x .
 \end{array}                                            \tag{D.50i}
\]
The free universal property gives continuous
endomorphisms, and the displayed relations
$r^4=s^2=1$, $srs=r^{-1}$, verified on
$x,y$, give their inverses and make them
automorphisms.
For $c=[x,y]$ direct multiplication gives
\[
                   r(c)=x^{-1}cx,\qquad
                   s(c)=c^{-1}.                        \tag{D.50j}
\]
Thus their outer classes lie in $T$.
Their homology matrices are
\[
             R=\begin{pmatrix}0&-1\\1&0\end{pmatrix},
       \qquad S=\begin{pmatrix}0&1\\1&0\end{pmatrix},
                                                        \tag{D.50k}
\]
and the eight matrices $R^iS^j$
($0\leq i<4$, $0\leq j<2$) are distinct.
These maps therefore exhibit an embedded
$D_8$ in $T$. The subgroups
$\langle r\rangle$, $\langle r^2,s\rangle$,
$\langle s\rangle$, and the identity realize
$C_4$, $C_2\times C_2$, $C_2$, and $1$.

We finish with the asserted multiplier data.
D.43 makes every involution's multiplier a
global sign. A positive nonidentity involution
has determinant $1$ in (D.50c). Its order-two
matrix is diagonalizable in characteristic
zero, and faithfulness excludes $I$; its matrix
is consequently $-I$.
There is at most one such involution in any
finite $2$-subgroup.
For $C_2\times C_2$, simultaneous diagonalization
identifies its full image with the four sign
matrices above, giving exactly one positive
and two negative involutions.
For $D_8$, the square of a rotation is
positive by (D.50d), so every other involution
is negative. The reflections generate $D_8$,
and a rotation is a product of two reflections.
The homomorphic multiplier is therefore
globally $1$ on the rotations and globally
$-1$ on the reflections.

Finally let $a$ have order four in an arbitrary
cyclic subgroup. Its matrix in (D.50c) has
minimal polynomial $X^2+1$ and determinant one.
Thus $\lambda(a)_2=1$ and $\lambda(a^2)_2=1$.
D.43 applied to the involution $a^2$ gives
$\lambda(a^2)=1$ globally. The multiplier
homomorphism then gives (D.50b).
No assertion of conjugacy follows from this
matrix classification.

The finite certificate
\path{code/part_c/torus_two_subgroups_d50.py}
verifies the full-word maps (D.50i)--(D.50j),
their eight distinct matrices, and the unit
square computations modulo eight used in
(D.50h). The passage from that finite
calculation to all of $\Q_2$ is the valuation
argument just given.
\end{proof}
